% ================================================================
% PRIMARI
% ================================================================
\section{Processi Primari}

\subsection{Processo Di Fornitura}
\label{sec:fornitura}

\subsubsection{Descrizione}

Il processo di fornitura regola le attività con cui \textit{Synergo Unit} si pone
nella posizione di fornitore nei confronti del committente (Prof.\ Vardanega) e del
mentore/cliente (Zucchetti S.p.A.). Il Regolamento del Progetto Didattico stabilisce
le seguenti distinzioni di ruolo\textsubscript{G}:

\begin{itemize}
  \item il \textbf{proponente} è \textit{cliente} rispetto alle esigenze di prodotto
        e \textit{mentore} rispetto alle scelte di sviluppo;
  \item il \textbf{docente} è \textit{committente} di tutte le forniture;
  \item l'interazione fornitore-committente concerne vincoli contrattuali
        (scadenze, costi) e questioni regolamentari.
\end{itemize}

Le attività di valutazione dei capitolati\textsubscript{G}, selezione e candidatura sono state
completate nella fase CC. Nella fase RTB il processo regola le comunicazioni
formali con i referenti esterni, la produzione della documentazione richiesta
per le revisioni di avanzamento e la rendicontazione periodica al committente.

\subsubsection{Attività}

\paragraph{Valutazione Dei Capitolati}
Il gruppo ha analizzato i capitolati disponibili producendo il documento
\href{https://synergo-unit-unipd.github.io/Documentary-Repo/docs/CC/doc_gestionali/doc_esterni/valutazione_dei_capitolati/valutazione_dei_capitolati.pdf}
     {\underline{Valutazione Dei Capitolati}},
che illustra i criteri di confronto, i vantaggi e gli svantaggi di ciascun
capitolato, includendo il resoconto degli incontri esplorativi con i proponenti,
e motiva la scelta finale.

\paragraph{Dichiarazione Degli Impegni\textsubscript{G}}
A seguito della selezione del capitolato, il gruppo ha redatto la
\href{https://synergo-unit-unipd.github.io/Documentary-Repo/docs/CC/doc_gestionali/doc_esterni/dichiarazione_degli_impegni/dichiarazione_degli_impegni.pdf}
     {\underline{Dichiarazione Degli Impegni}},
che specifica:

\begin{itemize}
  \item scadenza ultima di consegna prevista;
  \item preventivo dei costi\textsubscript{G} finali calcolato secondo il
        regolamento (costo minimo ammesso: 12.000\,€ per gruppi di 7,
        proporzionalmente ridotto per gruppi di 6);
  \item totale delle ore produttive assegnate a persona, strettamente compreso
        nell'intervallo 80\,(min)~--~95\,(max) ore per membro;
  \item distribuzione delle ore per ruolo (non per persona) secondo i costi orari
        ufficiali riportati nella Sezione~\ref{sec:ruoli};
  \item rischi attesi e strategie di mitigazione;
  \item impegni formali assunti verso il committente.
\end{itemize}

Il costo accettato all'aggiudicazione costituisce il limite superiore invalicabile
per l'intero svolgimento del progetto.

\paragraph{Comunicazioni Esterne}
Le comunicazioni ufficiali con aziende e docenti avvengono tramite l'indirizzo email
istituzionale del gruppo. Tale indirizzo è configurato con inoltro automatico verso
gli indirizzi personali dei membri delegati; le risposte vengono inviate mantenendo
come mittente l'indirizzo ufficiale, garantendo continuità e tracciabilità delle
comunicazioni formali. 

Le candidature alle revisioni di avanzamento avvengono per
posta elettronica, senza invio di allegati, esponendo il puntatore alla sezione del
repository\textsubscript{G} dedicata alla revisione\textsubscript{G}.

\paragraph{Incontri Con Il Proponente}
Ogni incontro con il proponente viene pianificato con congruo
anticipo, prevede la predisposizione di un ordine del giorno e si conclude con la
redazione e approvazione di un verbale esterno secondo le norme descritte nella
Sezione~\ref{sec:documentazione}.

\paragraph{Diario Di Bordo\textsubscript{G}}
Il gruppo è tenuto a relazionare al committente negli eventi "Diario di Bordo"
secondo il calendario stabilito, esponendo il rapporto tra costi sostenuti e
avanzamento conseguito e comunicando le difficoltà riscontrate.

\subsection{Processo Di Sviluppo}
\label{sec:sviluppo}

\subsubsection{Descrizione}

Il processo di sviluppo è attivo dalla baseline RTB e regola le attività di
analisi dei requisiti e progettazione del sistema. Le attività di codifica
saranno attivate nella fase PB\textsubscript{G}.

\subsubsection{Modello Di Sviluppo}

Il gruppo adotta il framework\textsubscript{G} \textbf{Scrum} con sprint\textsubscript{G} della durata di
\textbf{una settimana}, con cadenza da martedì a martedì. I ruoli vengono ruotati
a ogni sprint secondo le regole descritte nella Sezione~\ref{sec:ruoli}.

La pianificazione dettagliata degli sprint, il piano di rotazione dei ruoli e le
milestone\textsubscript{G} sono definiti nel documento \link{https://synergo-unit-unipd.github.io/Documentary-Repo/docs/RTB/doc_gestionali/doc_esterni/piano_di_progetto/piano_di_progetto.pdf}{Piano di Progetto\textsubscript{G}}.

\subsubsection{Analisi Dei Requisiti}

L'\link{https://synergo-unit-unipd.github.io/Documentary-Repo/docs/RTB/doc_tecnici/doc_esterni/analisi_dei_requisiti/analisi_dei_requisiti.pdf}{Analisi dei Requisiti} è il documento tecnico principale della fase RTB. I casi
d'uso sono organizzati in tre macroaree funzionali:

\begin{itemize}
  \item \textbf{Macroarea 1 — Gestione Editor e Markdown\textsubscript{G}}: scrittura, formattazione\textsubscript{G},
        rendering.
  \item \textbf{Macroarea 2 — Interazione AI (LLM\textsubscript{G})}: riassunto, riscrittura,
        traduzione\textsubscript{G}, i "Sei cappelli per pensare\textsubscript{G}", distant writing\textsubscript{G}.
  \item \textbf{Macroarea 3 — Gestione I/O e Persistenza}: salvataggio, apertura,
        esportazione note.
\end{itemize}

Per ogni macroarea i requisiti opzionali sono raggruppati in una sezione dedicata.

\paragraph{Convenzioni Dei Casi D'uso\textsubscript{G}}

I casi d'uso adottano le seguenti convenzioni:

\begin{itemize}
      \item \textbf{Denominazione}: formato \texttt{UC X: Testo Del Titolo}, dove
            \texttt{X} è il numero progressivo.
      \item \textbf{Ordine dei diagrammi}: i casi d'uso, nel foglio condiviso Lucidspark\textsubscript{G}, sono ordinati numericamente
            dall'alto verso il basso; le sottoclassi da sinistra verso destra.
      \item \textbf{Strumento}: i diagrammi UML\textsubscript{G} sono gestiti tramite
            \textbf{LucidChart} su un foglio condiviso, mantenuto aggiornato
            in modo incrementale.
      \item \textbf{Link interni}: i riferimenti tra casi d'uso all'interno
            dell'Analisi dei Requisiti\textsubscript{G} utilizzano i comandi
            \verb|\hyperlink| e \verb|\hypertarget| con il colore primario del
            template\textsubscript{G} e testo sottolineato, garantendo navigabilità
            nel documento PDF. La configurazione è definita nel preambolo
            condiviso (\texttt{shared/}).
      \item \textbf{Inclusioni ed Estensioni}: nel caso d'uso padre si aggiunge nel campo
        dedicato:
      \begin{verbatim}
            \item \textbf{Inclusioni}: \hyperlink{ucY.y}{UC Y.y: Titolo}.
            \item \textbf{Estensioni}: \hyperlink{ucX.x}{UC X.x: Titolo}.
      \end{verbatim}
      Nel caso d'uso inclusione o estensione si inserisce l'ancora prima del
        \verb|\paragraph|:
      \begin{verbatim}
      \hypertarget{ucX.x}{}
      \paragraph{UC X.x: Titolo}
      \end{verbatim}
\end{itemize}

\paragraph{Tecnologie: Valutazione, Confronto E Scelta}
 
Nell'ambito del processo di sviluppo, la gestione delle tecnologie
segue una responsabilità divisa per fasi:
 
\begin{itemize}
  \item \textbf{Valutazione e confronto}: spettano all'\textbf{Analista\textsubscript{G}},
        che produce il documento \textit{Analisi Dello Stato Dell'Arte}
        con pro, contro e motivazioni per ciascuna alternativa esaminata.
 
  \item \textbf{Scelta finale}: spetta al \textbf{Progettista\textsubscript{G}},
        che seleziona le tecnologie sulla base dell'analisi prodotta
        dall'Analista e delle esigenze architetturali del sistema.
\end{itemize}
 
Lo stack tecnologico adottato per il PoC\textsubscript{G} è il seguente:
 
\begin{itemize}
  \item \textbf{Frontend}: Vue.js con CodeMirror come editor Markdown\textsubscript{G}.
  \item \textbf{Backend}: Python con framework FastAPI\textsubscript{G}.
  \item \textbf{Database}: PostgreSQL\textsubscript{G}.
  \item \textbf{Deployment}: Docker\textsubscript{G} (containerizzazione dei tre
        moduli frontend, backend, database).
\end{itemize}
 
La documentazione completa delle valutazioni comparative è riportata nel documento \link{https://synergo-unit-unipd.github.io/Documentary-Repo/docs/RTB/doc_gestionali/doc_interni/analisi_stato_arte/analisi_stato_arte.pdf}{Analisi Dello Stato Dell'Arte}.
 